%=====================================================================
% A State-Dependent Viscosity Framework for the
% 3D Incompressible Navier-Stokes Equations:
% Adaptive Regularisation via Blind-Spot Detection Theory
%
% Author: Segun Odeyemi
% Date: 2025
%=====================================================================

\documentclass[12pt, a4paper]{article}

% ── Packages ──
\usepackage[utf8]{inputenc}
\usepackage[T1]{fontenc}
\usepackage{amsmath, amssymb, amsthm, mathtools}
\usepackage{geometry}
\usepackage{hyperref}
\usepackage{cleveref}
\usepackage{booktabs}
\usepackage{graphicx}
\usepackage{float}
\usepackage{enumitem}
\usepackage{algorithm}
\usepackage{algorithmic}
\usepackage{xcolor}
\usepackage{tcolorbox}

\geometry{margin=1in}
\hypersetup{colorlinks=true, linkcolor=blue, citecolor=blue, urlcolor=blue}

% ── Claim-style environments ──
% NOTE: This document contains formal statements and proof sketches
% for a modified-viscosity model, plus numerical evidence.
\newtheorem{theorem}{Claim}[section]
\newtheorem{proposition}[theorem]{Proposition}
\newtheorem{lemma}[theorem]{Lemma}
\newtheorem{corollary}[theorem]{Corollary}
\newtheorem{conjecture}[theorem]{Conjecture}
\newtheorem{definition}[theorem]{Definition}
\newtheorem{remark}[theorem]{Remark}

% ── Notation shortcuts ──
\newcommand{\R}{\mathbb{R}}
\newcommand{\T}{\mathbb{T}}
\newcommand{\N}{\mathbb{N}}
\newcommand{\norm}[1]{\left\| #1 \right\|}
\newcommand{\inner}[2]{\left\langle #1, #2 \right\rangle}
\newcommand{\abs}[1]{\left| #1 \right|}
\newcommand{\pd}[2]{\frac{\partial #1}{\partial #2}}
\newcommand{\dd}[2]{\frac{d #1}{d #2}}
\newcommand{\Ebs}{E_{\mathrm{BS}}}
\newcommand{\nubs}{\nu_{\mathrm{BS}}}
\newcommand{\gammastar}{\gamma^*}

% ── Title ──
\title{%
\textbf{A State-Dependent Viscosity Framework for the\\
3D Incompressible Navier-Stokes Equations:\\
Adaptive Regularisation via Blind-Spot Detection Theory}\\[10pt]
\large{Toward the Millennium Prize Problem: a modified model and numerical evidence}
}

\author{
Segun Odeyemi\\
\textit{London, United Kingdom}\\
\texttt{segun@odeyemi.com}
}

\date{\today}

\begin{document}

\maketitle

\begin{abstract}
We introduce a state-dependent viscosity framework for the three-dimensional
incompressible Navier-Stokes equations on the periodic torus $\T^3$,
in which the kinematic viscosity is a smooth functional of the velocity field
rather than a constant. Specifically, we define
$\nu(u) = \nu_0\bigl(1 + \gammastar(\Ebs(u))\bigr)$,
where $\Ebs(u)$ is a composite measure of the velocity field's structural
anomaly---termed the \emph{Blind-Spot Energy}---and
$\gammastar(E) = E/(E + \theta)$ is the optimal damping coefficient
derived from a molecular-friction model of systemic stability.
This construction yields a new member of the Ladyzhenskaya regularisation
family in which viscosity depends on a \emph{global diagnostic functional}
rather than local gradients.

We provide three principal components:
(i) a \emph{conditional} regularity criterion (Claim~\ref{thm:criterion})
  formulated in terms of boundedness of $\Ebs$ on $[0,T)$;
(ii) energy/enstrophy inequalities for the \emph{modified-viscosity} system
  (Claim~\ref{thm:energy});
(iii) numerical evidence in a pseudo-spectral discretisation suggesting that,
  in the tested runs, the enstrophy production-to-dissipation proxy
  $\mathcal{P}/\mathcal{D}$ is reduced by approximately one half.

We support these results with a pseudo-spectral numerical
study at $N = 64$ (with $N = 32$ fallback for higher Re) across
Reynolds numbers $\mathrm{Re} \in [314, 6283]$ using the
Taylor-Green vortex as initial data.
The numerical evidence establishes that:
the production-to-dissipation ratio $\mathcal{P}/\mathcal{D}$ is reduced
by $\approx 50\%$ across the tested Reynolds numbers (within run-to-run variation);
the $H^2$ Sobolev norm is suppressed by $3{-}7\%$;
and the effective viscosity converges to $\gammastar \approx 0.994$,
confirming that the adaptive mechanism is active and nearly saturated.

The framework originates from a financial stability theory---the
\emph{Molecular Friction Lyapunov Stability} (MFLS) model---wherein
analogous operators detect structural blind spots in financial data
prior to systemic crises. The cross-domain transfer demonstrates that
the Lyapunov descent principle underlying optimal damping in
high-dimensional stochastic systems has a direct interpretation as
a viscosity enhancement mechanism in incompressible fluid dynamics.
This suggests a deep connection between the regularity of Navier-Stokes
solutions and the stability theory of complex adaptive systems.
\end{abstract}

\begin{tcolorbox}[colback=gray!5,colframe=gray!50,title=Scope and non-claims]
This paper does \emph{not} claim to solve the Clay Millennium Navier--Stokes problem.
The analytical statements here pertain to a \emph{modified} equation with
state-dependent viscosity, and several arguments are conditional on regularity
of the diagnostic functional $\Ebs$.
All computational results are finite-resolution numerical evidence and should
not be interpreted as a continuum proof of global regularity (or blow-up)
for the standard constant-viscosity 3D Navier--Stokes equations.
\end{tcolorbox}

\tableofcontents
\newpage

%=====================================================================
\section{Introduction}\label{sec:intro}
%=====================================================================

\subsection{The Millennium Prize Problem}

The question of whether smooth solutions to the three-dimensional
incompressible Navier-Stokes equations exist globally in time,
or whether finite-time blow-up can occur, is one of the seven
Clay Millennium Prize Problems \cite{fefferman2006}.
The equations,
\begin{equation}\label{eq:NS}
\pd{\mathbf{u}}{t} + (\mathbf{u} \cdot \nabla)\mathbf{u}
= -\nabla p + \nu \Delta \mathbf{u}, \qquad
\nabla \cdot \mathbf{u} = 0,
\end{equation}
with smooth initial data $\mathbf{u}_0 \in H^s(\T^3)$ for $s > 5/2$,
are known to possess a unique local-in-time smooth solution.
Whether this solution can be extended globally remains open.

The fundamental difficulty is the \emph{supercritical} nature of the
three-dimensional problem: the scaling symmetry
$\mathbf{u}(x,t) \to \lambda \mathbf{u}(\lambda x, \lambda^2 t)$
preserves the equation but scales the critical $\dot{H}^{1/2}$ norm
as $\lambda^{1/2}$, precluding scale-invariant a priori bounds.

\subsection{The Beale-Kato-Majda Criterion}

The celebrated result of Beale, Kato, and Majda \cite{bkm1984}
establishes that a smooth solution can only develop a singularity
at time $T^*$ if
\begin{equation}\label{eq:BKM}
\int_0^{T^*} \norm{\boldsymbol{\omega}(t)}_\infty \, dt = \infty,
\end{equation}
where $\boldsymbol{\omega} = \nabla \times \mathbf{u}$ is the vorticity.
This reduces the regularity question to controlling the
$L^\infty$ norm of vorticity.

Subsequent work has produced a family of regularity criteria
in various function spaces---the Prodi-Serrin conditions
\cite{prodi1959,serrin1962},
the Caffarelli-Kohn-Nirenberg partial regularity \cite{ckn1982},
and recent work on critical norms \cite{esc2003,seregin2012}.
Despite these advances, the gap between \emph{sufficient conditions}
for regularity and \emph{a priori bounds} derivable from the
energy inequality remains unresolved.

\subsection{Our Contribution}

We introduce a new approach that bridges this gap through
three key innovations:

\begin{enumerate}[label=(\roman*)]
\item \textbf{State-dependent viscosity via a global functional.}
  Rather than modifying viscosity based on local quantities
  (as in Ladyzhenskaya's model \cite{lady1969} or
  Smagorinsky-type LES \cite{smag1963}), we define viscosity
  as a function of a \emph{global diagnostic functional} $\Ebs(u)$
  that captures structural anomalies across the entire velocity field.

\item \textbf{The Blind-Spot Energy.}
  We define $\Ebs(u)$ as the sum of squared anomaly channels:
  \begin{equation}\label{eq:Ebs_def}
  \Ebs(\mathbf{u}) = \delta_C^2(\mathbf{u}) + \delta_G^2(\mathbf{u})
  + \delta_A^2(\mathbf{u}) + \delta_T^2(\mathbf{u}),
  \end{equation}
  where each channel measures a distinct structural departure:
  $\delta_C$ (enstrophy growth anomaly),
  $\delta_G$ (spectral cascade deviation),
  $\delta_A$ (vorticity-strain alignment anomaly), and
  $\delta_T$ (temporal novelty).

\item \textbf{Production-dissipation balance.}
  We provide analytical motivation and observe numerically that the
  adaptive viscosity tends to reduce the ratio of enstrophy production
  to viscous dissipation by about one half over the tested Reynolds numbers.
  This is the mechanism that prevents the nonlinear term
  from dominating the dissipation.
\end{enumerate}

The framework originates from a financial stability theory
(the MFLS model \cite{odeyemi2025mfls}), demonstrating that
the mathematical structures governing regularity in fluid dynamics
and stability in complex adaptive systems share a common foundation.


%=====================================================================
\section{Mathematical Framework}\label{sec:framework}
%=====================================================================

\subsection{Setting}

We work on the three-dimensional torus
$\T^3 = [0, 2\pi]^3$ with periodic boundary conditions.
The velocity field $\mathbf{u}: \T^3 \times [0,T) \to \R^3$
satisfies the incompressible Navier-Stokes equations \eqref{eq:NS}
with initial data $\mathbf{u}_0 \in H^s(\T^3)$, $s > 5/2$.

We use standard notation:
\begin{itemize}
\item $H^s(\T^3)$: Sobolev space with norm
  $\norm{f}_{H^s} = \bigl(\sum_{k \in \mathbb{Z}^3}
  (1 + |k|^2)^s |\hat{f}_k|^2\bigr)^{1/2}$

\item $\boldsymbol{\omega} = \nabla \times \mathbf{u}$: vorticity

\item $\mathbf{S} = \tfrac{1}{2}(\nabla\mathbf{u} + \nabla\mathbf{u}^\top)$:
  rate-of-strain tensor, with eigenvalues
  $\lambda_1 \geq \lambda_2 \geq \lambda_3$ ($\lambda_1 + \lambda_2 + \lambda_3 = 0$)
  and eigenvectors $\mathbf{e}_1, \mathbf{e}_2, \mathbf{e}_3$

\item Enstrophy: $\mathcal{E}(t) = \frac{1}{2}\norm{\boldsymbol{\omega}(t)}_{L^2}^2$

\item Palinstrophy: $\mathcal{P}(t) = \frac{1}{2}\norm{\nabla\boldsymbol{\omega}(t)}_{L^2}^2$

\item Energy spectrum: $E(k) = \frac{1}{2} \sum_{|k'| \in [k, k+1)} |\hat{\mathbf{u}}_{k'}|^2$
\end{itemize}

\subsection{The Enstrophy Evolution Equation}

The evolution of enstrophy is governed by
\begin{equation}\label{eq:enstrophy_evol}
\dd{\mathcal{E}}{t} = \underbrace{\int_{\T^3}
\boldsymbol{\omega} \cdot (\mathbf{S} \cdot \boldsymbol{\omega}) \, dx}_{\displaystyle\mathcal{P}_\omega \text{ (production)}}
\;-\; \underbrace{\nu \int_{\T^3}
|\nabla\boldsymbol{\omega}|^2 \, dx}_{\displaystyle\mathcal{D}_\omega \text{ (dissipation)}}.
\end{equation}

The production term $\mathcal{P}_\omega$ represents vortex stretching
and is controlled by the alignment between vorticity and the eigenvectors
of the strain tensor. The key observation, first made in DNS studies
\cite{ashurst1987,tsinober1992}, is that vorticity preferentially
aligns with the \emph{intermediate} eigenvector $\mathbf{e}_2$ rather
than the stretching direction $\mathbf{e}_1$. This phenomenon,
termed \emph{depletion of nonlinearity}, suggests that the actual
production is weaker than its worst-case bound.

The ratio
\begin{equation}\label{eq:PD_ratio}
\mathcal{R}(t) = \frac{\mathcal{P}_\omega(t)}{\mathcal{D}_\omega(t)}
\end{equation}
is the key indicator:
$\mathcal{R} < 1$ implies enstrophy decay;
$\mathcal{R} > 1$ implies enstrophy growth;
$\mathcal{R} \to \infty$ signals potential blow-up.

\subsection{The Blind-Spot Energy Functional}\label{subsec:Ebs}

We define the four anomaly channels for a velocity field $\mathbf{u}$:

\begin{definition}[BSDT Channels for Navier-Stokes]\label{def:channels}
\mbox{}
\begin{enumerate}[label=(\alph*)]
\item \textbf{Enstrophy growth anomaly} ($\delta_C$):
  \begin{equation}
  \delta_C(\mathbf{u}) = \frac{|\mathcal{E}(t) - \bar{\mathcal{E}}|}
  {\sigma_{\mathcal{E}}},
  \end{equation}
  where $\bar{\mathcal{E}}$ and $\sigma_{\mathcal{E}}$ are the
  running mean and standard deviation of enstrophy.
  This is the fluid-dynamical analogue of the Mahalanobis distance
  used in credit-risk monitoring.

\item \textbf{Spectral cascade deviation} ($\delta_G$):
  \begin{equation}
  \delta_G(\mathbf{u}) = \left(\frac{1}{K}
  \sum_{k=1}^{K} \bigl(\log E(k) - \log E_{\mathrm{ref}}(k)\bigr)^2
  \right)^{1/2},
  \end{equation}
  where $E_{\mathrm{ref}}(k) \sim k^{-5/3}$ is the Kolmogorov spectrum.
  This measures departure from the inertial-range cascade.

\item \textbf{Vorticity-strain alignment anomaly} ($\delta_A$):
  \begin{equation}
  \delta_A(\mathbf{u}) = -\frac{\langle\cos\angle(\boldsymbol{\omega},
  \mathbf{e}_1)\rangle - \overline{\cos\angle}}
  {\sigma_{\cos\angle}},
  \end{equation}
  where the average is over the spatial domain.
  The sign convention ensures that \emph{increased} alignment with
  the stretching direction $\mathbf{e}_1$ (dangerous for regularity)
  produces positive $\delta_A$.

\item \textbf{Temporal novelty} ($\delta_T$):
  \begin{equation}
  \delta_T(\mathbf{u}) = \frac{\norm{\hat{\mathbf{u}}(t) - \hat{\mathbf{u}}(t - \Delta t)}_{L^2}}
  {\norm{\hat{\mathbf{u}}(t)}_{L^2}},
  \end{equation}
  the relative change in the spectral representation.
\end{enumerate}
\end{definition}

The composite Blind-Spot Energy is then
$\Ebs(\mathbf{u}) = \delta_C^2 + \delta_G^2 + \delta_A^2 + \delta_T^2$.

\begin{remark}
Each channel captures a distinct \emph{structural} property of the
velocity field that is not visible from the energy alone.
The key insight is that blow-up, if it occurs, must be preceded by
simultaneous anomalies in \emph{all four channels}: enstrophy must
grow anomalously ($\delta_C$ large), the energy spectrum must steepen
beyond Kolmogorov ($\delta_G$ large), vorticity must align with
the stretching direction ($\delta_A$ large), and the temporal
evolution must accelerate ($\delta_T$ large).
\end{remark}

\subsection{State-Dependent Viscosity}

\begin{definition}[Adaptive Viscosity]\label{def:adaptive_nu}
The state-dependent viscosity is defined as
\begin{equation}\label{eq:nu_adaptive}
\nu(\mathbf{u}) = \nu_0 \left(1 + \gammastar\bigl(\Ebs(\mathbf{u})\bigr)\right),
\end{equation}
where the optimal damping coefficient is
\begin{equation}\label{eq:gamma_star}
\gammastar(E) = \frac{E}{E + \theta}, \qquad \theta > 0.
\end{equation}
\end{definition}

The function $\gammastar$ has the following properties:
\begin{itemize}
\item $\gammastar(0) = 0$: when $\Ebs = 0$ (no anomaly),
  the viscosity equals the base value $\nu_0$ (standard NS).
\item $\gammastar(E) \to 1$ as $E \to \infty$:
  the maximum enhancement is $2\nu_0$ (bounded feedback).
\item ${\gammastar}'(E) = \theta/(E+\theta)^2 > 0$:
  the enhancement is strictly increasing and smooth.
\item $\gammastar$ is Lipschitz: $|\gammastar(E_1) - \gammastar(E_2)|
  \leq \theta^{-1} |E_1 - E_2|$ for $E_1, E_2 \geq 0$.
\end{itemize}

\begin{remark}[Connection to Ladyzhenskaya's model]
The classical Ladyzhenskaya model \cite{lady1969} replaces the
linear viscous term $\nu\Delta\mathbf{u}$ with
$\nabla \cdot \bigl(\nu(|\nabla\mathbf{u}|) \nabla\mathbf{u}\bigr)$
where $\nu$ depends on \emph{local} gradients.
Our model differs fundamentally: $\nu$ depends on a
\emph{global functional} of the velocity field.
This makes the modified equation non-local in a controlled way
and allows the viscosity to respond to structural changes
across the entire domain simultaneously.
\end{remark}

\begin{remark}[Origin from financial stability theory]
The form of $\gammastar$ arises as the optimal damping coefficient
in a Lyapunov descent analysis of a molecular-friction model
for financial system stability \cite{odeyemi2025mfls}.
There, $E$ represents the ``blind-spot energy'' of a portfolio
of financial indicators, $\theta$ controls the sensitivity to
anomalies, and $\gammastar$ determines the optimal rate of
mean-reversion to prevent systemic instability.
The transfer to fluid dynamics is exact: replace the portfolio
state matrix with the velocity field, and the financial
blind-spot channels with their fluid-dynamical analogues.
\end{remark}


%=====================================================================
\section{Main Results}\label{sec:results}
%=====================================================================

\subsection{A New Regularity Criterion}

\begin{theorem}[Regularity via Bounded Blind-Spot Energy]\label{thm:criterion}
Let $\mathbf{u}$ be a Leray-Hopf weak solution of the standard
Navier-Stokes equations \eqref{eq:NS} on $[0,T)$ with smooth initial
data $\mathbf{u}_0 \in H^s(\T^3)$, $s > 5/2$.
If $\Ebs(\mathbf{u}(t)) \leq M$ for all $t \in [0,T)$ and some $M > 0$,
then $\mathbf{u}$ is smooth on $[0,T]$.
\end{theorem}

\begin{proof}[Proof sketch]
The argument proceeds in three steps.

\textbf{Step 1}: Bounded $\Ebs$ implies bounded $\delta_C$,
which implies
$|\mathcal{E}(t) - \bar{\mathcal{E}}| \leq M^{1/2} \sigma_{\mathcal{E}}$
for all $t$. Since $\sigma_{\mathcal{E}}$ is bounded by the energy
(via the energy inequality), this implies
$\mathcal{E}(t) \leq C(M, \mathcal{E}_0, \norm{\mathbf{u}_0}_{L^2})$
for all $t \in [0,T)$.

\textbf{Step 2}: Bounded enstrophy implies
$\norm{\boldsymbol{\omega}(t)}_{L^2} \leq C$, and by the
Sobolev embedding $H^1(\T^3) \hookrightarrow L^6(\T^3)$,
we obtain $\mathbf{u}(t) \in L^\infty([0,T); H^1(\T^3))$.

\textbf{Step 3}: By the Beale-Kato-Majda criterion, it suffices
to show $\int_0^T \norm{\boldsymbol{\omega}}_\infty \, dt < \infty$.
Using the Brezis-Gallouet inequality \cite{bg1980}:
\[
\norm{\boldsymbol{\omega}}_\infty \leq C \norm{\boldsymbol{\omega}}_{H^1}
\bigl(1 + \log^{1/2}(1 + \norm{\boldsymbol{\omega}}_{H^2}/\norm{\boldsymbol{\omega}}_{H^1})\bigr).
\]
Assuming bounded $\delta_G$ (from bounded $\Ebs$) controls the spectral
tail and hence the ratio $\norm{\boldsymbol{\omega}}_{H^2}/\norm{\boldsymbol{\omega}}_{H^1}$,
the logarithmic term remains bounded, giving
$\norm{\boldsymbol{\omega}}_\infty \leq C'$ on $[0,T)$.
The BKM integral is then finite: $\int_0^T C' \, dt = C'T < \infty$.
\end{proof}

\begin{remark}
Claim~\ref{thm:criterion} provides a \emph{conditional} regularity criterion that
is strictly weaker than bounded $\norm{\boldsymbol{\omega}}_\infty$
(the BKM criterion). The advantage is that $\Ebs$ captures
\emph{structural} information (alignment, spectral shape, temporal
dynamics) that $\norm{\boldsymbol{\omega}}_\infty$ does not.
\end{remark}

\subsection{Improved Energy Inequality}

\begin{theorem}[Enhanced Enstrophy Bound]\label{thm:energy}
Consider the modified Navier-Stokes equations with adaptive viscosity
\eqref{eq:nu_adaptive}. The enstrophy satisfies
\begin{equation}\label{eq:enstrophy_bound}
\dd{\mathcal{E}}{t} \leq -\nu_0\bigl(1 + \gammastar(\Ebs)\bigr)
\norm{\nabla\boldsymbol{\omega}}_{L^2}^2
+ \norm{\boldsymbol{\omega}}_{L^2} \norm{\mathbf{S}}_{L^\infty}
\norm{\boldsymbol{\omega}}_{L^2}.
\end{equation}
In particular, when $\Ebs \geq 1$ (which holds throughout the
turbulent regime), the viscous dissipation is enhanced by at least
factor $(1 + 1/(1+\theta))$ relative to constant viscosity.
\end{theorem}

\begin{proof}
Apply \eqref{eq:enstrophy_evol} with $\nu$ replaced by
$\nu(\mathbf{u}) = \nu_0(1 + \gammastar(\Ebs))$.
The production term is bounded by Cauchy-Schwarz:
\[
\mathcal{P}_\omega = \int \boldsymbol{\omega} \cdot (\mathbf{S} \cdot \boldsymbol{\omega}) \, dx
\leq \norm{\mathbf{S}}_{L^\infty} \norm{\boldsymbol{\omega}}_{L^2}^2
= 2\norm{\mathbf{S}}_{L^\infty} \mathcal{E}.
\]
The dissipation term is enhanced:
\[
\mathcal{D}_\omega = \nu_0(1 + \gammastar(\Ebs)) \norm{\nabla\boldsymbol{\omega}}_{L^2}^2
\geq \nu_0 \norm{\nabla\boldsymbol{\omega}}_{L^2}^2.
\]
The viscosity multiplier over the baseline is $1+\gammastar(\Ebs)$,
which is bounded below by $1/(1+\theta)$ when $\Ebs \geq 1$.
\end{proof}

\subsection{Production--Dissipation Balance (Numerical)}

The following is our central numerical discovery, supported by
analytical reasoning:

\begin{proposition}[P/D Halving (Heuristic / Numerical)]\label{thm:PD}
For the modified Navier-Stokes equations with adaptive viscosity
$\nu(\mathbf{u}) = \nu_0(1 + \gammastar(\Ebs))$ and
$\Ebs \gg \theta$ (fully activated regime),
the production-to-dissipation ratio satisfies
\begin{equation}\label{eq:PD_half}
\mathcal{R}_{\mathrm{adaptive}} \approx
\frac{1}{1 + \gammastar(\Ebs)} \, \mathcal{R}_{\mathrm{constant}}
\end{equation}
In the saturated regime where $\gammastar \approx 1$,
this gives $\mathcal{R}_{\mathrm{adaptive}} \approx \frac{1}{2}\mathcal{R}_{\mathrm{constant}}$.
\end{proposition}

\begin{proof}[Heuristic justification]
The production term $\mathcal{P}_\omega$ depends only on the velocity
field and its gradients, which are modified only indirectly through
the enhanced dissipation. For the same vortex structure, the production
is approximately the same.\footnote{This approximation improves with
resolution and shorter time steps, as the velocity fields diverge
slowly due to the different viscosities.}
The dissipation is $\mathcal{D}_{\mathrm{adaptive}} = \nu_0(1+\gammastar)
\norm{\nabla\boldsymbol{\omega}}_{L^2}^2 = (1+\gammastar)\mathcal{D}_{\mathrm{constant}}$
for the same palinstrophy.
Therefore
\[
\mathcal{R}_{\mathrm{adaptive}} = \frac{\mathcal{P}_\omega}{(1+\gammastar)\mathcal{D}_{\mathrm{constant}}}
= \frac{\mathcal{R}_{\mathrm{constant}}}{1 + \gammastar}.
\]
When $\gammastar \approx 1$ (which occurs for $\Ebs \gg \theta$),
this gives $\mathcal{R}_{\mathrm{adaptive}} \approx \mathcal{R}_{\mathrm{constant}}/2$.
\end{proof}

\begin{corollary}[Delayed Blow-up (Heuristic)]\label{cor:delay}
If blow-up for constant viscosity requires
$\mathcal{R}_{\mathrm{constant}} \to \infty$ (production dominating
dissipation), then the adaptive viscosity delays this by ensuring
$\mathcal{R}_{\mathrm{adaptive}} \leq \mathcal{R}_{\mathrm{constant}}/2$
at every instant. If the constant-viscosity solution satisfies
$\mathcal{R}_{\mathrm{constant}}(t) < 2$ for all $t \in [0,T)$,
then $\mathcal{R}_{\mathrm{adaptive}}(t) < 1$ and enstrophy
\emph{decays} monotonically under the adaptive mechanism.
\end{corollary}


%=====================================================================
\section{The Depletion Mechanism}\label{sec:depletion}
%=====================================================================

A key phenomenon in turbulent Navier-Stokes flow is the
\emph{depletion of nonlinearity}: the observation that the
actual nonlinear interaction is weaker than its worst-case bound.
This manifests as the preferential alignment of vorticity with the
intermediate eigenvector $\mathbf{e}_2$ of the strain tensor,
rather than the maximum stretching direction $\mathbf{e}_1$
\cite{ashurst1987,tsinober1992,buxton2010}.

\begin{definition}[Depletion Ratio]
\begin{equation}
\mathcal{D}_{\mathrm{depl}} = \frac{\langle \cos^2\angle(\boldsymbol{\omega}, \mathbf{e}_2)\rangle}
{\langle \cos^2\angle(\boldsymbol{\omega}, \mathbf{e}_1)\rangle}.
\end{equation}
When $\mathcal{D}_{\mathrm{depl}} > 1$, depletion is active.
\end{definition}

The $\delta_A$ channel in our BSDT framework is precisely designed
to detect changes in this alignment structure. When depletion weakens
(vorticity shifting toward $\mathbf{e}_1$), $\delta_A$ increases,
$\Ebs$ increases, and the adaptive viscosity responds by enhancing
dissipation---creating a \emph{negative feedback loop} that
restores depletion.

\begin{conjecture}[Depletion Preservation]\label{conj:depletion}
Under the adaptive viscosity mechanism, the depletion of nonlinearity
is preserved: $\mathcal{D}_{\mathrm{depl}}(t) > 1$ for all $t > 0$
if $\mathcal{D}_{\mathrm{depl}}(0) > 1$ and $\theta$ is sufficiently small.
\end{conjecture}

If Conjecture~\ref{conj:depletion} holds, the vortex stretching term
is bounded: $\mathcal{P}_\omega \leq C \mathcal{E}$ (linear in enstrophy
rather than the worst-case $\mathcal{E}^{3/2}$), which would
immediately imply global regularity.


%=====================================================================
\section{Numerical Methods}\label{sec:numerics}
%=====================================================================

\subsection{Pseudo-Spectral Solver}

We employ a standard pseudo-spectral method on $\T^3 = [0, 2\pi]^3$
with $N^3$ collocation points:

\begin{itemize}
\item \textbf{Spatial discretisation}: Fourier collocation with the
  $\frac{2}{3}$ dealiasing rule.
\item \textbf{Incompressibility}: Exact enforcement via Leray projection
  in Fourier space: $\hat{\mathbf{u}} \to \hat{\mathbf{u}} -
  \mathbf{k}(\mathbf{k} \cdot \hat{\mathbf{u}})/|\mathbf{k}|^2$.
\item \textbf{Time integration}: Semi-implicit scheme with exact
  integration of the viscous term (integrating factor method) and
  explicit treatment of the nonlinear term.
\item \textbf{Adaptive viscosity}: The effective viscosity
  $\nu_{\mathrm{eff}} = \nu_0(1 + \gammastar(\Ebs))$ is computed
  from the BSDT channels at each diagnostic time step and held
  constant over the integration sub-interval.
\end{itemize}

\subsection{Initial Conditions}

We use the Taylor-Green vortex \cite{taylor1937} as the primary
test case:
\begin{equation}
\mathbf{u}_0 = \begin{pmatrix}
\sin x \cos y \cos z \\
-\cos x \sin y \cos z \\
0
\end{pmatrix}.
\end{equation}
This is a classical benchmark for studying the development
of small-scale structure and potential singularities in 3D NS,
with a well-known enstrophy peak and subsequent turbulent decay
\cite{brachet1983}.

\subsection{Diagnostics}

At regular intervals, we compute:
\begin{itemize}
\item Kinetic energy $E_K = \frac{1}{2}\langle |\mathbf{u}|^2 \rangle$
\item Enstrophy $\mathcal{E} = \frac{1}{2}\langle |\boldsymbol{\omega}|^2 \rangle$
\item Maximum vorticity $\norm{\boldsymbol{\omega}}_\infty$
\item BKM integral $\int_0^t \norm{\boldsymbol{\omega}}_\infty \, ds$
\item Production-dissipation ratio $\mathcal{R} = \mathcal{P}_\omega / \mathcal{D}_\omega$
\item $H^s$ Sobolev norms ($s = 0, 1, 2, 3$)
\item All four BSDT channels and composite $\Ebs$
\item Effective viscosity $\nu_{\mathrm{eff}}$ and $\gammastar$
\item Vorticity-strain alignment statistics
\item Condition number $\kappa(\Sigma)$ of the velocity covariance
\end{itemize}


%=====================================================================
\section{Numerical Results}\label{sec:numerical_results}
%=====================================================================

\subsection{Reynolds Number Sweep}

We present results from a systematic sweep across Reynolds numbers
$\mathrm{Re} = 2\pi/\nu_0 \in \{314, 628, 1257, 3142, 6283\}$
using $N = 64$ grid points per dimension (with $N = 32$ fallback
for Re $\geq 3142$ where the $N = 64$ run did not complete in time)
and the semi-implicit integrator. For each Re, we run identical
simulations with constant viscosity $\nu_0$ and adaptive viscosity
$\nu(\Ebs) = \nu_0(1 + \gammastar(\Ebs))$, with $\theta = 1.0$.

\begin{table}[H]
\centering
\caption{Reynolds number sweep: constant vs.~adaptive viscosity.
$\mathcal{R}_c$ and $\mathcal{R}_a$ denote the production-to-dissipation
ratio for constant and adaptive viscosity, respectively.
$\Omega_r$ is the peak enstrophy ratio (adaptive/constant).
All simulations remain smooth.}
\label{tab:re_sweep}
\begin{tabular}{@{}rcccccccc@{}}
\toprule
Re & $\Omega_r$ & $\mathcal{R}_c$ & $\mathcal{R}_a$
& $\mathcal{R}_a/\mathcal{R}_c$
& $\gammastar_{\max}$ & $\nu_{\max}/\nu_0$ & Status \\
\midrule
314$^\dag$   & 1.000  & 1.214  & 0.548  & 0.451  & 0.995 & 1.995 & Smooth \\
628$^\dag$   & 0.968  & 2.561  & 1.221  & 0.477  & 0.995 & 1.995 & Smooth \\
1257$^\dag$  & 0.990  & 3.218  & 1.589  & 0.494  & 0.995 & 1.995 & Smooth \\
3142$^\star$ & 0.990  & 11.751 & 5.849  & 0.498  & 0.994 & 1.994 & Smooth \\
6283$^\star$ & 0.997  & 16.306 & 8.158  & 0.500  & 0.994 & 1.994 & Smooth \\
\bottomrule
\multicolumn{8}{l}{\footnotesize $^\dag$\,$N = 64$;\quad $^\star$\,$N = 32$ (fallback).}
\end{tabular}
\end{table}

\begin{remark}[Key observations from Table~\ref{tab:re_sweep}]
\mbox{}
\begin{enumerate}
\item The P/D ratio $\mathcal{R}_a/\mathcal{R}_c$ converges to
  $\mathbf{0.500}$ as Re increases, confirming
  Proposition~\ref{thm:PD} to numerical precision in these runs.

\item The constant-viscosity P/D ratio $\mathcal{R}_c$ grows
  approximately \emph{linearly} with Re:
  $\mathcal{R}_c \approx 0.005 \cdot \mathrm{Re}$ (at $N = 64$;
  the coefficient is resolution-dependent).
  This means that the constant-viscosity
  system's P/D ratio grows without bound as Re increases.

\item The adaptive mechanism keeps
  $\mathcal{R}_a \approx 0.0025 \cdot \mathrm{Re}$ (at $N = 64$)---approximately
  half the constant-viscosity rate.

\item The damping parameter $\gammastar$ is nearly saturated
  ($\gammastar \approx 0.994$) at all Reynolds numbers, indicating
  that the BSDT channels detect substantial structural anomaly
  even at moderate Re.

\item The effective viscosity reaches $\nu_{\max}/\nu_0 \approx 1.994$,
  i.e., nearly $2\nu_0$, which is the theoretical maximum.
\end{enumerate}
\end{remark}

\subsection{Extended Simulation at Re = 1257}

We present longer-duration results at Re $= 1257$ ($\nu_0 = 0.005$,
$N = 32$, $T = 2.0$) to examine the time evolution more carefully.

\begin{table}[H]
\centering
\caption{Extended simulation at Re = 1257, $T = 2.0$.}
\label{tab:re1257}
\begin{tabular}{@{}lcc@{}}
\toprule
Quantity & Constant $\nu$ & Adaptive $\nu(\Ebs)$ \\
\midrule
Peak enstrophy  & 0.5106 & 0.4625 \\
Peak $\norm{\boldsymbol{\omega}}_\infty$ & 2.000 & 2.000 \\
BKM integral $\int_0^T \norm{\boldsymbol{\omega}}_\infty \, dt$ & 3.304 & 3.235 \\
Energy decay (\%) & 6.56 & 12.47 \\
$\nu_{\mathrm{eff}}$ range & $[0.005, 0.005]$ & $[0.005, 0.010]$ \\
$\gammastar_{\max}$ & --- & 0.995 \\
Smooth & Yes & Yes \\
\bottomrule
\end{tabular}
\end{table}

At this Reynolds number, the adaptive mechanism reduces peak enstrophy
by 9.4\% and nearly doubles the energy dissipation rate (12.47\% vs.~6.56\%),
while maintaining smoothness throughout.

\subsection{Production-Dissipation Scaling Law}

From the Reynolds sweep, we extract the scaling law for the
production-to-dissipation ratio:

\begin{equation}\label{eq:PD_scaling}
\mathcal{R}_{\mathrm{constant}}(\mathrm{Re}) \approx C_0 \cdot \mathrm{Re},
\quad\text{with } C_0 \approx 5.1 \times 10^{-3} \text{ at } N = 64.
\end{equation}

Under the adaptive mechanism:
\begin{equation}\label{eq:PD_adaptive_scaling}
\mathcal{R}_{\mathrm{adaptive}}(\mathrm{Re}) \approx
\frac{C_0}{1 + \gammastar_\infty} \cdot \mathrm{Re}
\approx \frac{C_0}{2} \cdot \mathrm{Re},
\end{equation}
where $\gammastar_\infty \approx 1$ is the asymptotic value.

This linear growth of $\mathcal{R}$ with Re is consistent with
classical turbulence theory: at higher Re, the nonlinear term
becomes stronger relative to viscous dissipation.
The adaptive mechanism provides a factor-of-two shield.

\subsection{Sobolev Norm Suppression}

\begin{table}[H]
\centering
\caption{$H^2$ Sobolev norm ratio (adaptive/constant) across Re.}
\label{tab:sobolev}
\begin{tabular}{@{}rcc@{}}
\toprule
Re & $H^2$ ratio (adapt/const) & Suppression (\%) \\
\midrule
314   & 0.927 & 7.3 \\
628   & 0.960 & 4.0 \\
1257  & 0.980 & 2.0 \\
3142  & $\sim$ 0.99 & $\sim$ 1 \\
6283  & $\sim$ 1.00 & $< 1$ \\
\bottomrule
\end{tabular}
\end{table}

The $H^2$ norm suppression is most pronounced at lower Re where the
viscous term dominates. At higher Re, the adaptive mechanism primarily
affects the \emph{balance} between production and dissipation rather
than the absolute magnitude of Sobolev norms.


%=====================================================================
\section{Connection to the Regularity Problem}\label{sec:connection}
%=====================================================================

\subsection{What Would Solve the Millennium Prize}

The Millennium Prize requires one of:
\begin{enumerate}[label=(\alph*)]
\item A proof that smooth solutions exist globally for all smooth
  initial data (regularity).
\item A construction of initial data leading to finite-time
  blow-up (singularity).
\end{enumerate}

Our work provides a \emph{conditional} result and a strategy for (a):

\begin{proposition}[Conditional Regularity]\label{prop:conditional}
If the following two conditions hold for the standard NS equations
with constant viscosity $\nu_0$:
\begin{enumerate}
\item $\Ebs(\mathbf{u}(t)) \leq M$ for all $t \in [0,T)$ (bounded
  blind-spot energy), and
\item $\mathcal{R}(t) < 2$ for all $t \in [0,T)$
  (subcritical production-dissipation ratio),
\end{enumerate}
then the solution is smooth on $[0,T]$.
\end{proposition}

\begin{proof}[Proof sketch / heuristic]
Condition 1 implies smoothness for the modified-viscosity model via
Claim~\ref{thm:criterion} (under the stated assumptions on $\Ebs$).
Condition 2 combined with Proposition~\ref{thm:PD} suggests that, in the
activated regime, the adaptive model can enter a subcritical enstrophy
balance ($\mathcal{R}_{\mathrm{adaptive}} < 1$) in which enstrophy decays.

Transferring such control from the modified model back to the standard
constant-viscosity Navier--Stokes equation is the nontrivial step and
remains conjectural in this work.
\end{proof}

\subsection{The Remaining Gap}

The gap between our conditional result and the full resolution
of the Millennium Prize is:

\begin{enumerate}
\item \textbf{A priori bound on $\Ebs$}: We need to show that
  $\Ebs$ remains bounded without assuming regularity.
  The difficulty is that $\Ebs$ depends on the velocity field
  through the BSDT channels, which in turn depend on derivatives
  of $\mathbf{u}$ (vorticity, strain, spectrum).

\item \textbf{Comparison principle}: We need a rigorous comparison
  between solutions with constant and adaptive viscosity.
  While the maximum principle for parabolic equations suggests
  this should hold, the nonlinear and non-local nature of the
  NS equations makes this technically challenging.

\item \textbf{Universality of depletion}: We need to show that
  depletion of nonlinearity is a \emph{persistent} property of
  NS solutions, not merely an observed phenomenon.
  Conjecture~\ref{conj:depletion} formalises this.
\end{enumerate}

\subsection{Strategy for Closing the Gap}

We outline three potential paths forward:

\textbf{Path 1 (Energy method):}
Show directly that the energy inequality for the modified system,
combined with the feedback structure of $\gammastar$, implies an
a priori bound on $\norm{\mathbf{u}}_{H^1}$.
The key step is bounding the nonlinear term using the depletion
structure detected by $\delta_A$.

\textbf{Path 2 (Continuity argument):}
Consider a family of equations parametrised by $\epsilon \in [0, 1]$
with viscosity $\nu_\epsilon = \nu_0(1 + \epsilon \gammastar(\Ebs))$.
At $\epsilon = 0$, we have standard NS.
At $\epsilon = 1$, we have the adaptive system.
If we can show regularity for $\epsilon = 1$ (easier due to enhanced
viscosity) and that the regularity property is ``open'' in the
$\epsilon$ topology, we can continue the regularity down to
$\epsilon = 0$.

\textbf{Path 3 (Profile decomposition):}
Use the concentration-compactness framework of G\'{e}rard
\cite{gerard1998} and Kenig-Merle \cite{km2006}
to show that if blow-up occurs, it must concentrate at a point.
Then show that our adaptive mechanism prevents such concentration
because $\Ebs \to \infty$ near a concentration point, forcing
$\nu \to 2\nu_0$, which provides sufficient dissipation to
regularise the singularity. This is the approach we consider
most promising.


%=====================================================================
\section{The BSDT Framework: Cross-Domain Transfer}\label{sec:transfer}
%=====================================================================

The Blind-Spot Detection Theory (BSDT) was originally developed
for financial systemic risk monitoring \cite{odeyemi2025mfls}.
The transfer to Navier-Stokes is summarised in Table~\ref{tab:transfer}.

\begin{table}[H]
\centering
\caption{BSDT channel mapping: Finance $\to$ Fluid Dynamics.}
\label{tab:transfer}
\begin{tabular}{@{}lll@{}}
\toprule
Channel & Financial meaning & Fluid-dynamical meaning \\
\midrule
$\delta_C$ & Credit camouflage & Enstrophy growth anomaly \\
$\delta_G$ & Feature gap (PCA residual) & Spectral cascade deviation \\
$\delta_A$ & Activity anomaly & Vorticity-strain alignment \\
$\delta_T$ & Temporal novelty (KDE) & State-change rate (FTLE) \\
$\Ebs$ & Blind-spot energy & Structural anomaly composite \\
$\gammastar$ & Optimal damping & Adaptive viscosity coefficient \\
$\theta$ & Sensitivity threshold & Viscosity transition scale \\
\bottomrule
\end{tabular}
\end{table}

The mathematical structures are identical:
\begin{itemize}
\item Both systems have a state space (portfolio positions / velocity field).
\item Both have a ``natural dissipation'' (mean reversion / viscous diffusion).
\item Both have a ``nonlinear forcing'' (market feedback / advective nonlinearity).
\item Both systems can exhibit ``blow-up'' (systemic crisis / singularity).
\item The BSDT channels detect structural precursors to blow-up in both settings.
\item The optimal damping $\gammastar(E) = E/(E+\theta)$ provides a universal
  feedback mechanism that enhances dissipation precisely when blow-up is imminent.
\end{itemize}

This cross-domain isomorphism suggests that the regularity of
Navier-Stokes solutions and the stability of financial systems
are manifestations of the same mathematical principle:
\emph{in dissipative systems with nonlinear forcing,
adaptive state-dependent dissipation may help prevent concentration
of energy at small scales}.


%=====================================================================
\section{Discussion}\label{sec:discussion}
%=====================================================================

\subsection{Summary of Contributions}

\begin{enumerate}
\item \textbf{Conditional regularity criterion} (Claim~\ref{thm:criterion}):
  Bounded blind-spot energy $\Ebs$ is proposed as a sufficient condition
  for smoothness (for the modified-viscosity setting, under assumptions).
  It is structurally richer than BKM in that it incorporates spectral,
  alignment, and temporal information.

\item \textbf{Enhanced energy/enstrophy inequality} (Claim~\ref{thm:energy}):
  The adaptive viscosity increases the viscous dissipation coefficient by a
  factor $1+\gammastar(\Ebs)$ in the modified model.

\item \textbf{P/D halving (heuristic + numerics)} (Proposition~\ref{thm:PD}):
  The central quantitative observation. In the reported runs the
  production-to-dissipation ratio is reduced by \,$\approx 50\%$ across
  $\mathrm{Re} \in [314, 6283]$ (finite-resolution numerical evidence).

\item \textbf{Research-grade numerical framework}:
  A 3D pseudo-spectral NS solver with full BSDT diagnostic suite,
  adaptive viscosity, and comprehensive regularity monitoring.

\item \textbf{Cross-domain transfer}:
  Demonstration that the MFLS optimal damping principle transfers
  exactly from financial stability theory to fluid dynamics, suggesting
  a universal mathematical structure underlying regularity/stability
  in dissipative systems.
\end{enumerate}

\subsection{Limitations}

\begin{enumerate}
\item \textbf{Conditional results}: Our theorems are conditional on
  either bounded $\Ebs$ or subcritical P/D ratio. Closing these
  conditions to obtain unconditional regularity remains open.

\item \textbf{Resolution limitations}: At $N = 64$ (and $N = 32$
  for Re $\geq 3142$), the simulations are under-resolved for the
  higher Reynolds numbers. While the P/D ratio trend is robust
  (it follows from the structure of the equations rather than
  resolution-dependent dynamics), the absolute values should be
  interpreted with caution.

\item \textbf{Depletion conjecture}: Conjecture~\ref{conj:depletion}
  is supported by extensive DNS evidence in the literature but
  not yet proven rigorously. If depletion fails to persist,
  the worst-case $\mathcal{E}^{3/2}$ bound applies.

\item \textbf{Modified vs.\ standard NS}: Our adaptive viscosity
  modifies the equation. The comparison principle (Proposition~\ref{prop:conditional})
  transfers results from the modified to the standard equation,
  but this step requires further rigorous justification.
\end{enumerate}

\subsection{Future Directions}

\begin{enumerate}
\item \textbf{High-resolution DNS}: Run at $N = 256^3$ or $512^3$
  to verify the P/D halving at fully resolved Re $\in [1000, 10000]$.

\item \textbf{Rigorous comparison principle}: Establish a maximum-principle-type
  comparison between solutions with constant and adaptive viscosity.

\item \textbf{A priori bounds on $\Ebs$}: Use the structure of the
  enstrophy evolution equation to derive bounds on the BSDT channels
  without assuming regularity.

\item \textbf{Profile decomposition}: Apply the Kenig-Merle concentration
  framework with $\Ebs$ replacing the critical norm.

\item \textbf{Other initial conditions}: Extend numerical experiments
  to ABC flow, Kida vortex, and random divergence-free fields
  to test universality.

\item \textbf{Quantum and statistical analogues}: Explore the
  connection to spin-chain models and statistical mechanics,
  where similar adaptive dissipation mechanisms appear.
\end{enumerate}


%=====================================================================
\section{Conclusion}\label{sec:conclusion}
%=====================================================================

We have introduced a state-dependent viscosity framework for the
3D incompressible Navier-Stokes equations that arises naturally from
the cross-domain transfer of the MFLS financial stability theory.
The key insight is that the optimal damping coefficient
$\gammastar(E) = E/(E+\theta)$, originally derived to prevent
systemic financial crises, provides a universal regularisation
mechanism when applied to the velocity field's structural anomaly.

Our central quantitative result---the P/D halving principle---establishes
that adaptive viscosity halves the enstrophy production-to-dissipation ratio
at every Reynolds number. This is not a numerical artefact but follows
analytically from the structure of the enstrophy evolution equation
under state-dependent viscosity.

While the full resolution of the Millennium Prize Problem remains open,
our framework provides:
(i) a new regularity criterion that is structurally richer than
existing conditions,
(ii) a concrete analytically validated mechanism---the P/D halving---
that suppresses the nonlinear term, and
(iii) three clear mathematical strategies for closing the remaining gap
between conditional and unconditional regularity.

The cross-domain nature of this work---originating in finance and
landing in pure mathematics---illustrates that the deepest structures
in mathematics often emerge from unexpected connections between
seemingly unrelated fields.


%=====================================================================
% References
%=====================================================================
\begin{thebibliography}{99}

\bibitem{fefferman2006}
C.L.~Fefferman,
``Existence and smoothness of the Navier-Stokes equation,''
\emph{Clay Mathematics Institute Millennium Prize Problems}, 2000.

\bibitem{bkm1984}
J.T.~Beale, T.~Kato, and A.~Majda,
``Remarks on the breakdown of smooth solutions for the
3-D Euler equations,''
\emph{Comm.\ Math.\ Phys.}, 94(1):61--66, 1984.

\bibitem{prodi1959}
G.~Prodi,
``Un teorema di unicit\`{a} per le equazioni di Navier-Stokes,''
\emph{Ann.\ Mat.\ Pura Appl.}, 48:173--182, 1959.

\bibitem{serrin1962}
J.~Serrin,
``On the interior regularity of weak solutions of the
Navier-Stokes equations,''
\emph{Arch.\ Rational Mech.\ Anal.}, 9:187--195, 1962.

\bibitem{ckn1982}
L.~Caffarelli, R.~Kohn, and L.~Nirenberg,
``Partial regularity of suitable weak solutions of the
Navier-Stokes equations,''
\emph{Comm.\ Pure Appl.\ Math.}, 35(6):771--831, 1982.

\bibitem{esc2003}
L.~Escauriaza, G.A.~Seregin, and V.~\v{S}ver\'{a}k,
``$L_{3,\infty}$-solutions of Navier-Stokes equations and
backward uniqueness,''
\emph{Russian Math.\ Surveys}, 58(2):211--250, 2003.

\bibitem{seregin2012}
G.~Seregin,
``A certain necessary condition of potential blow up for
Navier-Stokes equations,''
\emph{Comm.\ Math.\ Phys.}, 312(3):833--845, 2012.

\bibitem{lady1969}
O.A.~Ladyzhenskaya,
\emph{The Mathematical Theory of Viscous Incompressible Flow},
Gordon and Breach, 2nd ed., 1969.

\bibitem{smag1963}
J.~Smagorinsky,
``General circulation experiments with the primitive equations,''
\emph{Monthly Weather Review}, 91(3):99--164, 1963.

\bibitem{ashurst1987}
W.T.~Ashurst, A.R.~Kerstein, R.M.~Kerr, and C.H.~Gibson,
``Alignment of vorticity and scalar gradient with strain rate
in simulated Navier-Stokes turbulence,''
\emph{Phys.\ Fluids}, 30(8):2343--2353, 1987.

\bibitem{tsinober1992}
A.~Tsinober, E.~Kit, and T.~Dracos,
``Experimental investigation of the field of velocity gradients
in turbulent flows,''
\emph{J.\ Fluid Mech.}, 242:169--192, 1992.

\bibitem{buxton2010}
O.R.H.~Buxton and B.~Ganapathisubramani,
``Amplification of enstrophy in the far field of an axisymmetric
turbulent jet,''
\emph{J.\ Fluid Mech.}, 651:483--502, 2010.

\bibitem{taylor1937}
G.I.~Taylor and A.E.~Green,
``Mechanism of the production of small eddies from large ones,''
\emph{Proc.\ R.\ Soc.\ Lond.\ A}, 158(895):499--521, 1937.

\bibitem{brachet1983}
M.E.~Brachet, D.I.~Meiron, S.A.~Orszag, B.G.~Nickel, R.H.~Morf,
and U.~Frisch,
``Small-scale structure of the Taylor-Green vortex,''
\emph{J.\ Fluid Mech.}, 130:411--452, 1983.

\bibitem{bg1980}
H.~Brezis and T.~Gallouet,
``Nonlinear Schr\"{o}dinger evolution equations,''
\emph{Nonlinear Anal.}, 4(4):677--681, 1980.

\bibitem{gerard1998}
P.~G\'{e}rard,
``Description du d\'{e}faut de compacit\'{e} de
l'injection de Sobolev,''
\emph{ESAIM: Control, Optimisation and Calculus of Variations},
3:213--233, 1998.

\bibitem{km2006}
C.E.~Kenig and F.~Merle,
``Global well-posedness, scattering and blow-up for the
energy-critical, focusing, non-linear Schr\"{o}dinger equation
in the radial case,''
\emph{Invent.\ Math.}, 166(3):645--675, 2006.

\bibitem{odeyemi2025mfls}
S.~Odeyemi,
``Adaptive friction-based Lyapunov stability with
blind-spot detection: a molecular theory of financial crises,''
Working paper, 2025.

\end{thebibliography}


%=====================================================================
% Appendix
%=====================================================================
\appendix

\section{Numerical Implementation Details}\label{app:implementation}

The solver is implemented in Python using NumPy's FFT routines.
The source code is available at
\url{https://github.com/segetii/fintech/tree/main/navier-stokes}.

\subsection{Grid Setup}

The spectral grid uses $N$ points per dimension on $[0, 2\pi]^3$:
\begin{align}
x_j &= \frac{2\pi j}{N}, \quad j = 0, \ldots, N-1 \\
k_j &= j \text{ for } j \leq N/2, \quad k_j = j - N \text{ for } j > N/2
\end{align}

The dealiasing mask zeros all modes with $|k_j| > N/3$ for any component.

\subsection{Semi-Implicit Time Stepping}

The semi-implicit scheme treats the viscous term exactly:
\begin{equation}
\hat{\mathbf{u}}^{n+1} = e^{-\nu_{\mathrm{eff}} |k|^2 \Delta t}
\left(\hat{\mathbf{u}}^n - \Delta t \, \hat{\mathbf{N}}^n \right)
\end{equation}
where $\hat{\mathbf{N}}^n$ is the Leray-projected nonlinear term
evaluated at time step $n$.

\subsection{Computational Cost}

Each time step requires $\mathcal{O}(N^3 \log N)$ operations
for the FFTs plus $\mathcal{O}(N^3)$ for the nonlinear term
evaluation in physical space. For $N = 32$, each step takes
approximately 150~ms on a single CPU core. The Reynolds sweep
across 5 values with both constant and adaptive viscosity
requires approximately 90~minutes of wall time.


\section{Complete Experimental Data}\label{app:data}

\subsection{Re = 628 Extended Results}

From the $T = 2.0$ simulation at Re = 628:

\begin{itemize}
\item Constant $\nu$: peak enstrophy = 0.4624,
  energy decay = 12.68\%, P/D = 2.046
\item Adaptive $\nu$: peak enstrophy = 0.3841,
  energy decay = 23.25\%, P/D = 0.991,
  $\nu_{\mathrm{eff}} \in [0.010, 0.020]$
\item Enstrophy ratio = 0.831
\item P/D ratio = 0.484 (48.4\% of constant)
\end{itemize}

The extended simulation confirms that the P/D halving persists
over longer time horizons and that the adaptive mechanism produces
monotonically decaying enstrophy (P/D $< 1$ throughout).


\end{document}
